\chapter{Theorie der fraktalen Entropie}
\label{ch:fraktale_entropie}

\section{Einleitung}
\label{sec:entropie_einleitung}

Die Theorie der fraktalen Entropie ist eine fundamentale Erweiterung der WDBT+. Sie beschreibt einen geschlossenen, stationären Kreislauf der fraktalen Entropie zwischen Vakuum, Materie und dem \( \Omega \)-Feld. Die Gravitation erzeugt durch die DSTT-Drift (\( \dot{\beta} > 0 \)) irreversible Prozesse und einen intrinsischen Zeitpfeil, während die Elektrodynamik (\( \dot{\beta} = 0 \)) stabile Strukturen bewahrt.

Die fraktale Raumdimension \( D \approx 2,704 \) (siehe Kapitel~\ref{ch:fraktale_dimension}) verhindert ein globales Entropiemaximum und damit den Wärmetod. Energie, Masse und Dichte bleiben konstant; die fraktale Entropie zirkuliert zwischen den Phasen.

\section{Axiome der Theorie}
\label{sec:entropie_axiome}

\subsection{Axiom 1: Fraktale Raumstruktur}
\label{sec:entropie_axiom1}

Der fundamentale Raum ist eine fraktale Menge \( \mathcal{M} \) mit Hausdorff-Dimension \( D \):

\[
D = \frac{\ln(20\phi)}{\ln(2 + \phi)} \approx 2,704, \quad \phi = \frac{1 + \sqrt{5}}{2}
\]

Die fundamentale Längenskala ist \( l_0 \approx 1,8 \times 10^{-15} \, \text{m} \) (siehe Kapitel~\ref{ch:bec_objekte}).

\subsection{Axiom 2: Fraktale Entropie}
\label{sec:entropie_axiom2}

Für eine Wahrscheinlichkeitsdichte \( \rho : \mathcal{M} \to \mathbb{R}^+ \) mit \( \int_{\mathcal{M}} \rho \, d^D r = 1 \) ist die fraktale Entropie definiert als:

\[
S_D[\rho] = -k_B \int_{\mathcal{M}} d^D r \, \rho(\vec{r}) \ln \rho(\vec{r}) \cdot \left( \frac{|\vec{r}|}{l_0} \right)^{D-3}
\]

Für \( D = 3 \) ergibt sich die Boltzmann-Entropie. Für \( D < 3 \) ist \( S_D \) subextensiv: Sie wächst langsamer als das Volumen.

\subsection{Axiom 3: Globaler Entropieerhaltungssatz}
\label{sec:entropie_axiom3}

Die Gesamt-fraktale Entropie des Universums ist konstant:

\[
\frac{dS_D^{\text{ges}}}{dt} = 0
\]

\subsection{Axiom 4: Lokale Entropieproduktion}
\label{sec:entropie_axiom4}

In jeder einzelnen Phase gilt der zweite Hauptsatz in erweiterter Form:

\[
\frac{dS_D^{(i)}}{dt} \geq 0
\]

\subsection{Axiom 5: Gravitation als Quelle der Irreversibilität}
\label{sec:entropie_axiom5}

Die Weber-Gravitation mit \( \beta_{\text{WG}} = 0,5 \) (siehe Kapitel~\ref{ch:dstt_weber}) führt zu einer monotonen Zunahme des tangentialen Energieanteils:

\[
\dot{\beta}_{\text{DSTT}} > 0 \quad \text{für alle } t
\]

\subsection{Axiom 6: Elektrodynamik als Stabilisator}
\label{sec:entropie_axiom6}

Die Weber-Elektrodynamik mit \( \beta_{\text{WED}} = 2 \) führt zu keiner Änderung der Energieaufteilung:

\[
\dot{\beta}_{\text{DSTT}} = 0 \quad \text{für alle } t
\]

\subsection{Axiom 7: Beschränktheit aller physikalischen Größen}
\label{sec:entropie_axiom7}

Für jede physikalische Größe \( X \) (Energie, Masse, Dichte, Temperatur, Volumen, Komplexität) gilt:

\[
\limsup_{t \to \infty} |X(t)| < \infty
\]

Keine physikalische Größe divergiert oder wächst unbegrenzt.

\section{Die drei Phasen}
\label{sec:entropie_phasen}

Das Universum besteht aus drei Phasen, zwischen denen die fraktale Entropie zirkuliert:

\begin{itemize}
    \item \textbf{Vakuum (Nullpunkt):} \( S_D^{\text{vac}} \) – die fraktale Struktur des Raums selbst
    \item \textbf{Materie (kinetisch):} \( S_D^{\text{kin}} \) – Bewegung der Teilchen, unterteilt in:
        \begin{itemize}
            \item \( S_D^{\text{kin,grav}} \) – durch Gravitation gebundene Materie
            \item \( S_D^{\text{kin,em}} \) – durch Elektrodynamik gebundene Materie
        \end{itemize}
    \item \textbf{\( \Omega \)-Feld:} \( S_D^{\Omega} \) – das Gravitationsfeld als Rotationsfeld (siehe Kapitel~\ref{ch:maxwell_gravitation})
\end{itemize}

Der Erhaltungssatz lautet:

\[
S_D^{\text{vac}}(t) + S_D^{\text{kin,grav}}(t) + S_D^{\text{kin,em}}(t) + S_D^{\Omega}(t) = S_D^{\text{ges}} = \text{const}
\]

\section{Die Entropieflüsse}
\label{sec:entropie_fluesse}

Aus den Bewegungsgleichungen der WDBT+ folgen die Entropieflussraten:

\[
\begin{aligned}
\frac{dS_D^{\text{vac}}}{dt} &= -\sigma_{\text{cre}} + \sigma_{\text{back}} \\
\frac{dS_D^{\text{kin,grav}}}{dt} &= +\sigma_{\text{cre}} - \sigma_{\text{drift}} \\
\frac{dS_D^{\text{kin,em}}}{dt} &= 0 \quad \text{(elektromagnetische Prozesse sind reversibel)} \\
\frac{dS_D^{\Omega}}{dt} &= +\sigma_{\text{drift}} - \sigma_{\text{back}}
\end{aligned}
\]

mit:

\begin{itemize}
    \item \( \sigma_{\text{cre}} > 0 \): Entropiefluss von Vakuum zu gravitativer Materie (Materieentstehung durch Quantenpotential, siehe Anhang~\ref{app:energiebilanz})
    \item \( \sigma_{\text{drift}} > 0 \): Entropiefluss von gravitativer Materie zum \( \Omega \)-Feld (DSTT-Drift)
    \item \( \sigma_{\text{back}} > 0 \): Entropiefluss vom \( \Omega \)-Feld zum Vakuum (Rückkopplung)
\end{itemize}

\section{Der geschlossene Kreislauf}
\label{sec:entropie_kreislauf}

Der geschlossene Entropiekreislauf lässt sich wie folgt darstellen:

\[
\text{Vakuum} \; \xrightarrow{\sigma_{\text{cre}}} \; \text{Materie (grav.)} \; \xrightarrow{\sigma_{\text{drift}}} \; \Omega\text{-Feld} \; \xrightarrow{\sigma_{\text{back}}} \; \text{Vakuum}
\]

Die elektromagnetische Materie ist vom Kreislauf entkoppelt (reversibel).

\subsection{Stationärer Zustand}
\label{sec:entropie_stationaer}

Im stationären Zustand gilt:

\[
\sigma_{\text{cre}} = \sigma_{\text{drift}} = \sigma_{\text{back}} \equiv \sigma
\]

Damit sind die Entropien jeder Phase konstant:

\[
\frac{dS_D^{\text{vac}}}{dt} = \frac{dS_D^{\text{kin,grav}}}{dt} = \frac{dS_D^{\text{kin,em}}}{dt} = \frac{dS_D^{\Omega}}{dt} = 0
\]

Es zirkuliert nur der Fluss, nicht die Menge.

\section{Warum kein Wärmetod eintritt}
\label{sec:entropie_waermetod}

\subsection{Bedingungen für einen Wärmetod}
\label{sec:entropie_waermetod_bedingungen}

Ein Wärmetod tritt ein, wenn:

\begin{enumerate}
    \item Die Entropie ein globales Maximum erreicht
    \item Alle Energiegradienten verschwinden
    \item Keine irreversiblen Prozesse mehr stattfinden
    \item Das System in einen Gleichgewichtszustand übergeht
\end{enumerate}

\subsection{Die Verhinderung des Wärmetods}
\label{sec:entropie_verhinderung}

Die fraktale Entropie \( S_D \) hat bei fester Gesamtenergie kein globales Maximum. Die Zustandsdichte im fraktalen Raum skaliert mit \( E^{D/3} \). Für \( D < 3 \) ist die Funktion

\[
S_D(E) = \frac{D}{3}k_B \ln E + \text{const}
\]

monoton wachsend und konkav, besitzt aber kein endliches Maximum.

Die DSTT-Drift (\( \dot{\beta} > 0 \)) erzeugt permanent radiale Energiegradienten:

\[
\nabla E_{\text{kin,grav}} \neq 0 \quad \text{für alle } t
\]

Die Raten \( \sigma_{\text{cre}}, \sigma_{\text{drift}}, \sigma_{\text{back}} \) sind immer positiv. Das System hat keinen attraktiven Fixpunkt im Phasenraum, sondern einen attraktiven Grenzzyklus.

\subsection{Hauptsatz der fraktalen Thermodynamik}
\label{sec:entropie_hauptsatz}

Unter den Axiomen 1–7 gilt:

\begin{enumerate}
    \item \textbf{Lokale Entropieproduktion:} \( \frac{dS_D^{(i)}}{dt} \geq 0 \) für jede Phase
    \item \textbf{Globale Entropieerhaltung:} \( \frac{dS_D^{\text{ges}}}{dt} = 0 \)
    \item \textbf{Permanente Dynamik:} \( \sigma_{\text{cre}}, \sigma_{\text{drift}}, \sigma_{\text{back}} > 0 \) für alle \( t \)
    \item \textbf{Kein Entropiemaximum:} Es existiert kein endlicher Zustand maximaler Gesamtentropie
    \item \textbf{Kein Wärmetod:} Das System erreicht nie einen Gleichgewichtszustand mit verschwindenden Gradienten und Prozessen
\end{enumerate}

\section{Kosmologische Konsequenzen}
\label{sec:entropie_kosmologie}

Die Theorie der fraktalen Entropie führt zu einem stationären Universum (siehe auch Kapitel~\ref{ch:kosmologie}):

\begin{itemize}
    \item Der Raum expandiert nicht, aber Strukturen und Bahnen dehnen sich durch die DSTT-Drift aus.
    \item Die mittlere Dichte bleibt durch kontinuierliche Materieentstehung konstant (siehe Anhang~\ref{app:energiebilanz}).
    \item Die Gravitation produziert permanent Entropie – dies ist der intrinsische Zeitpfeil.
    \item Die Elektrodynamik ermöglicht reversible Prozesse und bewahrt stabile Strukturen.
    \item Der Wärmetod tritt nicht ein – das Universum bleibt ewig dynamisch.
\end{itemize}

\section{Zusammenfassung}
\label{sec:entropie_zusammenfassung}

Die Theorie der fraktalen Entropie liefert ein konsistentes, logisch geschlossenes Bild:

\begin{itemize}
    \item Der Raum ist fraktal mit \( D \approx 2,704 \) und besitzt eine fundamentale Länge \( l_0 \).
    \item Die fraktale Entropie \( S_D \) ist subextensiv und zirkuliert in einem geschlossenen Kreislauf zwischen Vakuum, Materie und \( \Omega \)-Feld.
    \item Die Gravitation (\( \beta_{\text{WG}} = 0,5 \)) erzeugt durch \( \dot{\beta} > 0 \) irreversible Prozesse und den intrinsischen Zeitpfeil.
    \item Die Elektrodynamik (\( \beta_{\text{WED}} = 2 \)) mit \( \dot{\beta} = 0 \) ermöglicht reversible Prozesse und bewahrt stabile Strukturen.
    \item Der Wärmetod wird durch die fraktale Geometrie (kein globales Entropiemaximum) und die permanenten irreversiblen Flüsse verhindert.
    \item Das Universum ist stationär: Energie, Masse und Dichte bleiben konstant; die Dynamik zirkuliert ewig.
\end{itemize}